% DSC 148 - Course Project Report  (OFFICIAL ACM TEMPLATE)
% Compile on Overleaf with the "acmart" class (ACM Conference Proceedings).
% Figures expected at: outputs/figX_*.png
\documentclass[sigconf,nonacm]{acmart}

\usepackage{booktabs}
\usepackage{graphicx}
\usepackage{amsmath}

\settopmatter{printacmref=false}
\renewcommand\footnotetextcopyrightpermission[1]{}
\setcopyright{none}
\pagestyle{plain}

\newcommand{\D}{\Delta}

\begin{document}

\title{Beyond the Snapshot: Trajectory-Aware Features for Win and
Comeback Prediction in Professional League of Legends}

\author{Evan Ngo}
\affiliation{%
  \institution{University of California, San Diego}
  \department{Hal\i c\i o\u{g}lu Data Science Institute}
  \city{La Jolla}\state{CA}\country{USA}}
\email{mr.evanngo@gmail.com}

\author{Jimbo Cai}
\affiliation{%
  \institution{University of California, San Diego}
  \department{Hal\i c\i o\u{g}lu Data Science Institute}
  \city{La Jolla}\state{CA}\country{USA}}
\email{cjimbo29@gmail.com}

\begin{abstract}
Win prediction in professional esports is usually set up as binary classification on game-state
features taken at a single fixed timestamp. That snapshot setup throws away something basic:
which direction the game is actually moving. We build a trajectory-aware feature set that adds
momentum terms to the usual 10- and 15-minute differentials, measuring how gold, experience (XP),
and kill leads change between the two checkpoints, and we test it across four model families
(Logistic Regression, Gaussian Naive Bayes, Random Forest, and LightGBM) on 5{,}000 professional
matches from the 2024 Oracle Elixir dataset across ten leagues. We also define a separate comeback
task: among teams down more than 500 gold at ten minutes, predict which ones come back and win. We
report what we found without dressing it up. After feature engineering, a well-regularized
Logistic Regression is the best model on both tasks (win AUC 0.785, comeback AUC 0.734), and
gradient boosting does not beat it. That is a useful negative result given how often this kind of
problem gets handed to a heavier model by default. An ablation shows that objective control, not
trajectory, is the signal that actually matters: dropping objective features costs 0.023 AUC,
while dropping trajectory features costs only 0.002. Even so, the single most-split feature in
LightGBM is gold momentum, which exposes a gap between how much a feature gets used and how much
unique information it carries.
\end{abstract}

\keywords{League of Legends, esports analytics, win prediction, trajectory features,
gradient boosting, comeback classification, ablation}

\maketitle

\section{Introduction}
Competitive esports has turned into a genuinely useful testbed for data mining, mostly because the
match logs are so detailed: you get fine-grained game state across millions of games. League of
Legends (LoL), made by Riot Games, is one of the most-played and most-watched titles in the world,
with a pro scene covering Korea (LCK), China (LPL), Europe (LEC), North America (LCS), and a long
list of regional leagues.

The core question in this space is easy to state: given the game state at some early, fixed point
in time, how likely is a team to win? Most existing work treats this as a static classification
problem, featurizing the state at 10 or 15 minutes and fitting a classifier on those numbers. The
issue is that this ignores where the game is heading. A team sitting on a 1{,}000-gold lead that
climbed from zero at ten minutes is in a totally different spot than a team on a 1{,}000-gold lead
that fell from 3{,}000. The first is snowballing; the second is bleeding out a lead. A snapshot
model cannot tell them apart.

This paper has three parts. \textbf{(1)} We add a trajectory-aware feature set on top of the usual
snapshot features, with momentum terms that track how gold, XP, and kill differentials move
between 10 and 15 minutes, and we test it across four model families instead of pinning it to one.
\textbf{(2)} We define a comeback task, predicting recovery for teams that are behind at ten
minutes, and treat it as its own harder, class-imbalanced problem. \textbf{(3)} We keep the
comparison honest. Instead of crowning a complicated model, we show that a regularized linear model
comes out on top once the features are engineered well, that trajectory information gets used a lot
but is not actually irreplaceable, and that objective control is what drives the predictions. We
think that is more useful to know than squeezing out a fraction of a point on a leaderboard.

\section{Dataset}
\subsection{Source and Scope}
We use the 2024 Oracle Elixir professional match dataset~\cite{oracleselixir}, which is the
standard public source for pro LoL data. It pulls match data from the Riot Games API across every
major and regional league at multiple in-game timestamps. The full dataset is roughly 60{,}000
team- and player-level rows across 5{,}000 unique professional matches. Our win-prediction work
happens at the team-game level, and we use the Blue-side view, which gives 5{,}000 labeled
instances. We could double that to 10{,}000 by adding Red-side rows, which we leave as an easy
extension.

\begin{figure*}[t]
  \centering
  \includegraphics[width=0.92\textwidth]{outputs/fig1_dataset_overview.png}
  \caption{Dataset overview for Oracle Elixir 2024. Top: games per league, win rate by side,
  games per month. Bottom: game-length distribution (mean 42.6 minutes), gold differential at
  15 minutes (Blue perspective), and Blue win rate by patch.}
  \label{fig:overview}
\end{figure*}

\subsection{Key Statistics}
The data cover all ten primary leagues. By match count the biggest are LPL (1{,}032 games,
20.6\%), LCK (1{,}011, 20.2\%), LEC (645, 12.9\%), and LCS (483, 9.7\%), with six regional circuits
making up the rest (Figure~\ref{fig:overview}). Matches run from January to December 2024 across
patches 14.1 through 14.16. Game lengths sit between 25 and 60 minutes with a mean of 42.6, and the
Blue-side win rate is 51.4\%, the familiar small Blue-side edge that shows up in our data. Once we filter to complete records, the key columns have no missing
values, so we did not need to impute anything.

\subsection{Exploratory Findings}
Three things from EDA shaped the design. First, gold differential at 15 minutes is the strongest
single linear predictor of the outcome ($r=0.40$), with kill differential at 15 minutes close
behind ($r=0.34$). Mid-game state carries more signal than early state, which is why we keep the
15-minute features. Second, the outcome depends heavily on the early gold state
(Figure~\ref{fig:comeback}): teams ahead at ten minutes win 71.2\% of the time, even games land
near 51.6\%, and teams down 500 or more gold win only 31.2\%. The comeback rate falls off fast as
the deficit grows, from 36.0\% at a 500 to 1{,}000 gold deficit down to 6.5\% past 2{,}000 gold,
which is what makes the comeback task both well-defined and actually hard. Third, objective counts
have clearly nonlinear relationships with win rate, which is the main reason we bring in tree-based
models alongside the linear ones.

\begin{figure}[t]
  \centering
  \includegraphics[width=\columnwidth]{outputs/fig3_comeback_analysis.png}
  \caption{Comeback analysis. Left: win rate by gold state at ten minutes. Center: gold
  differential at ten minutes by outcome. Right: comeback rate vs.\ deficit magnitude.}
  \label{fig:comeback}
\end{figure}

\section{Predictive Tasks}
\subsection{Task 1: Win Prediction}
The main task is binary classification: given the team-level features we can see through fifteen
minutes, predict whether the team wins. \textbf{Metrics.} We report Accuracy, F1, and AUC-ROC under
5-fold stratified cross-validation. AUC-ROC is the headline number because it does not depend on a
threshold and it is what most prior work reports, so it makes comparison easier. \textbf{Baselines.}
Logistic Regression (linear log-odds), Gaussian Naive Bayes (conditional independence), and Random
Forest (nonparametric ensemble) cover a range of modeling assumptions, and LightGBM is the
nonlinear model we measure against them.

\subsection{Task 2: Comeback Classification}
The second task only looks at games where Blue is down more than 500 gold at ten minutes
(1{,}307 games, 26.1\% of the Blue-side sample) and asks whether that team recovers. It is
imbalanced (31.2\% positive) and on purpose harder than Task 1. \textbf{Metrics.} F1 is the main
number given the imbalance, and we also report AUC-ROC, Precision, and Recall, with class-weight
balancing on every model.

\section{Methodology}
\subsection{Feature Engineering}
We build 21 features in six groups (Table~\ref{tab:features}). The trajectory group is the new
piece. Gold momentum is
\begin{equation}
  \text{gold\_momentum} = \text{GD}_{15} - \text{GD}_{10},
\end{equation}
which is positive when a lead is growing (snowballing) and negative when it is shrinking; XP and
kill momentum work the same way. The interaction terms capture joint effects that matter more
together than apart, for example
\begin{equation}
  \text{gold\_x\_baron} = \text{GD}_{15}\times(\text{barons}-\text{opp\_barons}),
\end{equation}
since a gold lead plus baron control is worth more than either one on its own.

\begin{table}[t]
\caption{Engineered feature groups (21 features).}
\label{tab:features}
\begin{tabular}{@{}llc@{}}
\toprule
Group & Features & \# \\
\midrule
Early (10m)  & GD, XPD, CSD, KD at 10m & 4 \\
Mid (15m)    & GD, XPD, CSD, KD at 15m & 4 \\
Trajectory   & $\D$GD, $\D$XPD, $\D$KD (10$\to$15m) & 3 \\
Objectives   & dragon/baron/herald/tower diff, plates, first-obj & 6 \\
Interactions & GD$_{15}\!\times\!$baron, GD$_{15}\!\times\!$dragon & 2 \\
Context      & patch (numeric), league win rate & 2 \\
\bottomrule
\end{tabular}
\end{table}

\subsection{Models}
We treat all four models as candidates for the same feature set rather than picking a favorite up
front. \textbf{Logistic Regression} (standardized inputs, $C{=}1.0$) is the linear reference.
\textbf{Naive Bayes} is the lower bound. \textbf{Random Forest} (200 trees, depth 8) is a strong
nonparametric baseline. \textbf{LightGBM}~\cite{lightgbm} is the gradient-boosted candidate
(500 estimators, learning rate 0.05, 31 leaves, depth 6, subsample and colsample 0.8,
$\ell_1/\ell_2{=}0.1$, 50-round early stopping).

\section{Related Work}
Win prediction in MOBA games goes back to at least 2014. Yang et al.~\cite{yang2014} modeled
in-game combat as a sequence of graphs and showed that patterns in those combat graphs predict
game outcomes. Semenov et al.~\cite{semenov2016} predicted the match winner from team drafts in
Dota~2, testing Naive Bayes, logistic regression, gradient-boosted trees, and factorization
machines (which worked best), and reported AUC in the 0.66 to 0.71 range depending on skill tier.
Our Task 1 sits a bit higher (0.785), which is expected since we use in-game state through fifteen
minutes rather than drafts alone. Hodge et al.~\cite{hodge2019} built the first live professional
win-prediction system deployed at a real tournament, using engineered features with standard ML
models rather than a new algorithm. Closer to our setting, Maymin~\cite{maymin2021} calibrated an
in-game win-probability model on millions of ranked LoL games, with a focus on valuing individual
player actions. We do
three things differently: we pull the trajectory between two checkpoints out as explicit,
interpretable features instead of learning it from a full sequence; we treat comeback as its own
task; and we run a controlled comparison whose main takeaway is a caution rather than a new
state-of-the-art number.

\section{Results}
\subsection{Task 1: Win Prediction}
Table~\ref{tab:task1} and Figure~\ref{fig:models} show the 5-fold CV results. The clearest finding
is that \textbf{Logistic Regression is the best model on every metric} (accuracy 0.711, F1 0.720,
AUC 0.785). LightGBM is right there (AUC 0.771) but never passes it, Random Forest is just behind
(0.770), and Naive Bayes is the weakest (0.754). Our read is that once you engineer the trajectory
and interaction terms by hand, what is left of the relationship between features and win
probability is close to linear, so there is not much nonlinear structure left for boosting to find.

\begin{table}[t]
\caption{Task 1 win prediction (5-fold stratified CV). Best per column in bold.}
\label{tab:task1}
\begin{tabular}{@{}lccc@{}}
\toprule
Model & Accuracy & F1 & AUC-ROC \\
\midrule
\textbf{Logistic Regression} & \textbf{0.711} & \textbf{0.720} & \textbf{0.785} \\
Naive Bayes   & 0.687 & 0.693 & 0.754 \\
Random Forest & 0.694 & 0.703 & 0.770 \\
LightGBM      & 0.700 & 0.710 & 0.771 \\
\bottomrule
\end{tabular}
\end{table}

\begin{figure*}[t]
  \centering
  \includegraphics[width=0.92\textwidth]{outputs/fig6_model_comparison.png}
  \caption{Task 1 model comparison with CV error bars. Top-three differences are small relative to
  fold-to-fold variance.}
  \label{fig:models}
\end{figure*}

\subsection{Task 2: Comeback Prediction}
Table~\ref{tab:task2} tells the same story. Logistic Regression leads on F1 (0.551) and AUC
(0.734). LightGBM gets a decent AUC (0.702) but the \emph{worst} F1 (0.406): its probability
ranking is fine, but its default 0.5 threshold is too cautious on the minority class, which kills
recall (0.384). Every model drops well below its Task 1 level, which confirms that predicting a
comeback is just harder than predicting outcomes overall. Calibrating the threshold (Platt scaling,
say) is the obvious next move.

\begin{table}[t]
\caption{Task 2 comeback prediction (5-fold stratified CV). Best per column in bold.}
\label{tab:task2}
\begin{tabular}{@{}lcccc@{}}
\toprule
Model & F1 & AUC-ROC & Prec & Recall \\
\midrule
\textbf{Logistic Regression} & \textbf{0.551} & \textbf{0.734} & 0.470 & \textbf{0.667} \\
Naive Bayes   & 0.546 & 0.717 & 0.472 & 0.649 \\
Random Forest & 0.516 & 0.712 & \textbf{0.523} & 0.510 \\
LightGBM      & 0.406 & 0.702 & 0.439 & 0.384 \\
\bottomrule
\end{tabular}
\end{table}

\subsection{Ablation Study}
We drop one feature group at a time from a lighter LightGBM (Figure~\ref{fig:ablation}). That
smaller config sits a bit below the tuned model in Table~\ref{tab:task1} in absolute AUC (0.753),
so the relative drops are the thing to read. \textbf{Dropping Objectives causes by far the biggest
hit} (0.023 AUC, down to 0.730), which says objective control is doing most of the work. Dropping
Trajectory costs only 0.002 AUC, small but consistent across folds. Oddly, dropping the Mid (15m)
group nudges AUC up (to 0.756), which points to redundancy between the 15-minute snapshot and the
trajectory and early features.

\begin{figure}[t]
  \centering
  \includegraphics[width=\columnwidth]{outputs/fig8_ablation.png}
  \caption{Ablation: AUC-ROC when each feature group is removed (dashed line = full model).}
  \label{fig:ablation}
\end{figure}

\subsection{Feature Importance}
By LightGBM split count (Figure~\ref{fig:importance}), the most-used feature is gold momentum,
ahead of XP differential at 15 minutes and gold differential at 10 minutes. The interesting part is
how this clashes with the ablation: trajectory is the feature the model splits on most, but pulling
it out barely moves AUC. The reason is that momentum is just a deterministic function of the 10-
and 15-minute differentials, so the information is still there even without it. Split count
measures how handy a feature is for the trees, not how much unique signal it adds.

\begin{figure}[t]
  \centering
  \includegraphics[width=\columnwidth]{outputs/fig9_feature_importance.png}
  \caption{Top-15 LightGBM features by split importance. Gold momentum ranks first despite small
  ablation impact.}
  \label{fig:importance}
\end{figure}

\subsection{Patch Sensitivity}
Training on patches 14.1 through 14.11 and testing on 14.12 through 14.16, LightGBM hits AUC 0.761,
close to its CV value of 0.771. So the feature set holds up across patches, which makes sense given
that objective-control mechanics stay fairly stable patch to patch in pro play.

\subsection{Case Study}
Take a game where Blue is up 500 gold at ten minutes but has a gold momentum of $-800$, so it is
sitting at $-300$ by fifteen minutes. A snapshot model calls that close to a coin flip. A
trajectory-aware model gives it a lower win probability, because the negative momentum is a sign
that Red is taking over, usually by trading the early lead for objectives. The interaction of
\texttt{kd10}, \texttt{gd15}, and \texttt{gold\_momentum} is what lets the model separate ``ahead
and fading'' from ``ahead and pulling away,'' which is exactly the thing a snapshot loses.

\section{Discussion}
\textbf{Why the linear model wins.} The feature engineering does a lot of the work a nonlinear
model would otherwise have to learn on its own. Momentum is a hand-computed difference, and the
interaction terms hand the model the products a tree would need a bunch of splits to build. What is
left is close to linearly separable. \textbf{Used vs.\ irreplaceable.} The gold-momentum result is
a good reminder not to read split-based importance as causal contribution; importance and ablation
are answering two different questions. \textbf{Limitations.} We only use team-level aggregates, so
player skill, champion matchups, and draft synergy are all unobserved. We model 5{,}000 Blue-side
games, and adding Red-side rows would double that. The comeback labels use a 500-gold cutoff at ten
minutes, and a different cutoff could shift the conclusions. \textbf{Broader use.} Win-probability
models are useful for broadcast overlays, coaching tools, and draft simulators, and the comeback
classifier is handy for flagging the high-leverage moments.

\section{Conclusion}
We built a trajectory-aware feature set for win and comeback prediction in pro League of Legends
and tested it across four model families on 5{,}000 Oracle Elixir 2024 matches. The honest version
of the results: a regularized Logistic Regression is the best model on both tasks, gradient
boosting adds nothing once the trajectory and interaction features exist, and objective control,
not trajectory, is the signal that matters most. Trajectory features still get used the most by the
boosted model, which comes down to redundancy with the raw differentials. The whole thing is
deployed as an interactive dashboard. The main lesson we took away is that in this setting, good
features with a simple model beat a complex model run on default settings.

\begin{acks}
Data provided by Oracle Elixir (\url{https://oracleselixir.com}).
\end{acks}

\begin{thebibliography}{9}
\bibitem{oracleselixir} Oracle Elixir. \emph{League of Legends Esports Match Data}, 2024.
  \url{https://oracleselixir.com}.
\bibitem{yang2014} P.~Yang, B.~Harrison, and D.~Roberts. Identifying patterns in combat that are
  predictive of success in MOBA games. \emph{Proc.\ Foundations of Digital Games}, 2014.
\bibitem{semenov2016} A.~Semenov, P.~Romov, S.~Korolev, D.~Yashkov, and K.~Neklyudov. Performance
  of machine learning algorithms in predicting game outcome from drafts in Dota~2.
  \emph{Analysis of Images, Social Networks and Texts (AIST)}, Springer CCIS, vol.~661, 2016.
\bibitem{hodge2019} V.~Hodge, S.~Devlin, N.~Sephton, F.~Block, P.~Cowling, and A.~Drachen. Win
  Prediction in Multi-Player Esports: Live Professional Match Prediction.
  \emph{IEEE Trans.\ Games}, 2019.
\bibitem{lightgbm} G.~Ke et al. LightGBM: A highly efficient gradient boosting decision tree.
  \emph{NeurIPS}, 30, 2017.
\bibitem{maymin2021} P.~Z.~Maymin. Smart kills and worthless deaths: eSports analytics for League
  of Legends. \emph{J.\ Quantitative Analysis in Sports}, 17(1):11--27, 2021.
\end{thebibliography}

\end{document}